\documentclass[a4paper,12pt]{article}

\usepackage{graphicx}
\usepackage{listings}
\usepackage{hyperref}

\title{Practical Work 1: TCP File Transfer}
\author{Do Ba Hoang Minh - 22BI13278}
\date{26/11/2024}

\begin{document}

\maketitle

\section*{Goal}
The objective of this project is to implement a file transfer system over TCP/IP using socket programming. The system includes:
\begin{itemize}
    \item A server that receives the file and saves it.
    \item A client that sends the file to the server.
    \item Communication via TCP sockets.
\end{itemize}

\section*{Protocol Design}
The protocol utilizes TCP (Transmission Control Protocol), which ensures reliable and ordered communication. The flow is as follows:
\begin{itemize}
    \item The server initializes a socket, binds to a port, and listens for incoming connections.
    \item The client connects to the server and sends file data in chunks.
    \item The server receives the data and writes it to a file on disk.
    \item Both client and server close the connection upon completion.
\end{itemize}

\begin{figure}[h!]
    \centering
    \includegraphics[width=0.8\textwidth]{tcp_flow_diagram.png} % Replace with your actual diagram.
    \caption{TCP File Transfer Protocol Flow}
\end{figure}

\section*{System Organization}
The system consists of two components:
\begin{itemize}
    \item **Server**: A Python script that listens for connections and receives files.
    \item **Client**: A Python script that reads a file and sends it to the server.
\end{itemize}

\begin{figure}[h!]
    \centering
    \includegraphics[width=0.8\textwidth]{system_organization.png} % Replace with your actual diagram.
    \caption{System Organization Diagram}
\end{figure}

\section*{Implementation}
The implementation is done in Python, with separate scripts for the server and client.

\subsection*{Server Code}
Below is the server implementation:
\begin{lstlisting}[language=Python, caption=Server Code: tcp_file_server.py]
import socket

def start_server(host='127.0.0.1', port=65432, output_file='received_file.txt'):
    server_socket = socket.socket(socket.AF_INET, socket.SOCK_STREAM)
    server_socket.bind((host, port))
    server_socket.listen(1)

    print(f"Server listening on {host}:{port}...")

    conn, addr = server_socket.accept()
    print(f"Connected by {addr}")

    with open(output_file, 'wb') as file:
        while True:
            data = conn.recv(1024)
            if not data:  # Connection closed
                break
            file.write(data)

    print(f"File received and saved as {output_file}.")
    conn.close()
    server_socket.close()

if __name__ == "__main__":
    start_server()
\end{lstlisting}

\subsection*{Client Code}
Below is the client implementation:
\begin{lstlisting}[language=Python, caption=Client Code: tcp_file_client.py]
import socket

def send_file(file_path, host='127.0.0.1', port=65432):
    client_socket = socket.socket(socket.AF_INET, socket.SOCK_STREAM)
    client_socket.connect((host, port))

    with open(file_path, 'rb') as file:
        while (chunk := file.read(1024)):
            client_socket.sendall(chunk)

    print(f"File {file_path} has been sent to the server.")
    client_socket.close()

if __name__ == "__main__":
    send_file('file_to_send.txt')  # Replace with the actual file you want to send
\end{lstlisting}

\section*{Conclusion}
This project successfully demonstrated file transfer over TCP using Python. It showcased the reliable nature of TCP and provided practical experience in socket programming.

\end{document}
